\begin{tikzpicture}[x=1cm,y=1cm]\path[use as bounding box] (0,0) rectangle (16.6,15.8);
\node[head] at (0.1,15.45){(a) Temperature truth and a lagged record};\draw[wire] (0.1,15.17)--(7.699999999999999,15.17);
\node[head] at (8.35,15.45){(b) The corrected signal can imitate health drift};\draw[wire] (8.35,15.17)--(16.4,15.17);
\begin{axis}[mathaxis,at={(.85cm,10.55cm)},anchor=south west,width=7.15cm,height=4.55cm,
 xmin=-.5,xmax=4,ymin=-.5,ymax=9,xtick={0,1,2,3,4},ytick={0,4,8},xlabel={$t/\tau$},ylabel={Temperature rise [K]},
 legend style={at={(.97,.06)},anchor=south east}]
\addplot[line width=.85pt] coordinates {(-.5,0)(0,0)(0,8)(4,8)};\addlegendentry{Junction truth}
\addplot[dashed,domain=0:4,samples=100,line width=.85pt]{8*(1-exp(-x))};\addlegendentry{Recorded indication}
\draw[<->] (axis cs:1,5.056964)--(axis cs:1,8) node[midway,left,tinytext] {$-e(t)$};\end{axis}
\node[eq] at (3.9,9.45){$\tau\dot T^{\mathrm{record}}+T^{\mathrm{record}}=T^{\mathrm{true}}$};
\begin{axis}[mathaxis,at={(9.05cm,10.55cm)},anchor=south west,width=7.15cm,height=4.55cm,
 xmin=0,xmax=4,ymin=0,ymax=27,xtick={0,1,2,3,4},ytick={0,5,10,20,25},xlabel={$t/\tau$},ylabel={Corrected change [mV]},
 legend style={at={(.97,.97)},anchor=north east}]
\addplot[domain=0:4,samples=100,line width=.85pt]{20*exp(-x)};\addlegendentry{Zero health change}
\addplot[domain=0:4,samples=100,dashed,line width=.85pt]{5+20*exp(-x)};\addlegendentry{5 mV health change}
\addplot[domain=0:4,samples=2,densely dotted,forget plot]{5};
\draw[densely dotted] (axis cs:1.386294,0)--(axis cs:1.386294,5);
\addplot[only marks,mark=o,mark size=2.2pt,forget plot] coordinates {(1.386294,5)};\end{axis}
\node[tinytext,align=center] at (12.35,9.45){Analytic step: $\Delta T=8$ K, $\alpha=2.5$ mV/K.};
\node[head] at (0.1,9.0){(c) Derive the sign and the settling budget};\draw[wire] (0.1,8.72)--(7.699999999999999,8.72);
\node[head] at (8.35,9.0){(d) Visual error budget for an estimable target};\draw[wire] (8.35,8.72)--(16.4,8.72);
\node[eq] at (3.9,8.1){$e(t)=T^{\mathrm{record}}-T^{\mathrm{true}}=-\Delta T e^{-t/\tau}$};
\node[eq] at (3.9,6.97){$\widehat{\Delta h}(t)=\Delta h-\alpha[e(t)-e_b]$\\[5pt]
$e_b=0\Rightarrow\widehat{\Delta h}(t)=\Delta h+20e^{-t/\tau}\ \mathrm{mV}$};
\node[eq] at (3.9,5.42){$|\alpha\Delta T|e^{-t/\tau}\leq B_T$\\[5pt]
$t\geq\tau\max\{0,\log(|\alpha\Delta T|/B_T)\}$};
\node[eq] at (3.9,4.26){$B_T=5\ \mathrm{mV}\Rightarrow t\geq\tau\log4\simeq1.386\tau$};
\node[tinytext,align=center] at (3.9,2.86){The threshold controls this thermal residual only.\\It does not by itself control random false alarms.\\The step and first-order lag are illustrative;\\a real lag bound needs sensor or device evidence.};
\draw[arr] (8.8,5.55)--(16.0,5.55) node[below] {$\Delta h$};
\draw[densely dotted,wire] (12.4,5.55)--(12.4,6.0);
\draw[line width=1.2pt] (10.6,6.0)--(14.2,6.0);
\draw[wire] (10.6,5.83)--(10.6,6.17);\draw[wire] (14.2,5.83)--(14.2,6.17);
\fill (12.4,6) circle(2pt);\node[above] at (12.4,6.18){$\widehat{\Delta h}$};
\draw[<->] (10.6,6.85)--(14.2,6.85) node[midway,above] {$2zs_\Delta$};
\draw[dashed,line width=.8pt] (9.6,6)--(10.6,6);\draw[dashed,line width=.8pt] (14.2,6)--(15.2,6);
\draw[wire] (9.6,5.7)--(9.6,6.3);\draw[wire] (15.2,5.7)--(15.2,6.3);
\draw[<->] (14.2,5.2)--(15.2,5.2) node[midway,below] {$B_T$};
\node[eq] at (12.35,8.06){$\widehat{\Delta h}=\Delta h+b_T+s_\Delta Z,\quad Z\sim\mathcal N(0,1)$};
\node[eq] at (12.35,4.01){$|b_T|\leq B_T,\quad z=\Phi^{-1}(1-\alpha_0/2)$\\[5pt]
$\mathcal I=[\widehat{\Delta h}-(zs_\Delta+B_T),\ \widehat{\Delta h}+(zs_\Delta+B_T)]$};
\node[tinytext,align=center] at (12.35,2.72){For known $s_\Delta$ and a valid deterministic bias bound,\\$\Pr(\Delta h\in\mathcal I)\geq1-\alpha_0$.\\Solid: Gaussian radius; dashed: bounded-bias extension.};
\draw[wire] (.1,1.35)--(16.4,1.35);
\node[tinytext,align=center] at (8.25,.72){Temperature correction $\longrightarrow$ residual lag $\longrightarrow$ target bias $\longrightarrow$ uncertainty allocation.\\Averaging windows and time-varying temperature require the complete estimator's weighted discrepancy operator.};
\end{tikzpicture}
